\begin{abstract}
Hierarchical federated learning keeps Internet-of-Things data local, but its
differential-privacy guarantee protects single records, whereas clinical
deployments must protect patients with many records. This paper proposes
a client--edge--cloud framework whose privacy unit is the patient. Each client
clips a patient's summed gradient, calibrates noise to its population and
accounts for every hierarchical release. The guarantee is proved for the training
transcript and, with public encoders and standardization constants, end to
end (eICU). It is evaluated on VitalDB, $187$
eICU hospitals and, at record level, a wearable corpus. At a
training budget of at most $\varepsilon = 0.9$ ($1.0$ including the
conformal release; $0.766$ at record level on WEAR), it reduces uplink traffic
by $90.8$--$96.1\,\%$ and modeled energy per round by $72.3$--$79.6\,\%$;
$95$th-percentile workstation-CPU inference latency is at most $2.3$~ms.
On person-partitioned VitalDB the guarantee costs no utility (area under the
receiver operating characteristic curve, AUC, $0.832$ vs.\ $0.820$
non-private). On eICU it costs $3.1$ points (inconclusive against the
two-point target) and $3.9$ end to end, missing it; a non-private
hierarchical baseline is more accurate ($0.825$ vs.\ $0.794$). It uploads $2.2$--$4.0$ times less than re-implemented hierarchical and private
baselines and leaks less in the VitalDB audit. On VitalDB and WEAR,
pretrained initialization, not head-only noise, makes a strict budget
survivable; the patient-unit cost follows $A(1 - 1/k)$ in records per patient
$k$. A membership-inference audit with a positive control finds record-level
training above its difficulty floor (attack AUC $0.576$) and patient-level
training at it ($0.503$).
\end{abstract}
